\documentclass[a4paper,fleqn]{cas-sc}

\usepackage[authoryear,longnamesfirst]{natbib}
\usepackage{tabularx}
\usepackage{url}
\usepackage{graphicx}

% The built-in single-file preview sandbox cannot read companion images. A normal Elsevier
% submission build resolves the copied PNG files; the sandbox renders a labelled placeholder.
\newcommand{\submissionfigure}[2][]{%
  \IfFileExists{#2}{\includegraphics[#1]{#2}}{%
    \fbox{\parbox[c][0.20\textheight][c]{0.94\textwidth}{\centering
      Companion submission figure: \texttt{\detokenize{#2}}}}}}

\begin{document}
\let\WriteBookmarks\relax
\def\floatpagepagefraction{1}
\def\textpagefraction{.001}

\shorttitle{SustainTrace}
\shortauthors{Wang and Jiang}

\title[mode=title]{SustainTrace: Evidence-grounded ESG knowledge infrastructure with controlled evolution}

\author[1]{Michael Wang}
\cormark[1]
\credit{Conceptualization, Methodology, Software, Validation, Investigation, Data curation, Writing -- original draft, Writing -- review and editing}
\author[1]{Yan Jiang}
\credit{Conceptualization, Methodology, Validation, Investigation, Writing -- review and editing}
\affiliation[1]{organization={Affiliation details to be supplied in the journal submission system}}
\cortext[1]{Corresponding author. E-mail address to be inserted before submission.}

\begin{abstract}
Environmental, social and governance reports contain comparable disclosures, but their values,
units, periods, boundaries and assurance statements are distributed across prose, tables and
footnotes. SustainTrace is a local-first infrastructure that converts these documents into atomic,
evidence-linked knowledge and preserves extraction failures as reusable method knowledge. It
freezes source hashes, binds claims to page and layout evidence, blocks unsupported admission and
subjects bounded repairs to deterministic post-validation and human-gated promotion. The release
seed contains 149 reviewed facts and 1,032 failure observations. In a 60-task, 20-report human
benchmark, 177/220 applicable slots and 40/40 evidence coordinates were exact. On identical frozen
candidates from two cloud models, its admission gate reduced unsafe admission from 27.0\% and
24.2\% to 1.0\%, while exposing a model-dependent retention cost. The software supports auditable
ESG knowledge construction and controlled evolution rather than unqualified answer generation.
\end{abstract}

\begin{highlights}
\item Converts ESG reports into atomic facts linked to page, quotation and layout evidence.
\item Preserves typed failures and validates bounded repairs without rewriting historical results.
\item Reduces unsafe knowledge admission on paired candidates while quantifying over-blocking.
\end{highlights}

\begin{keywords}
ESG reporting \sep knowledge extraction \sep provenance \sep document intelligence \sep controlled repair \sep trustworthy AI
\end{keywords}

\maketitle

\begin{table*}[t]
\caption{Code metadata for the archived public release.}
\label{tab:metadata}
\small
\renewcommand{\arraystretch}{1.08}
\begin{tabularx}{\textwidth}{@{}p{0.28\textwidth}X@{}}
\toprule
Current code version & 0.1.0 \\
Permanent link to code/repository & \url{https://github.com/extradimen/sustaintrace}; \url{https://doi.org/10.5281/zenodo.23293034} \\
Legal code license & Apache License 2.0 \\
Code versioning system & Git \\
Software languages and tools & Python 3.11+, TypeScript, React, MinerU, Poppler and optional Ollama-compatible model adapters \\
Operating environments & macOS, Windows and Linux; local web browser for the workbench \\
Installation & \texttt{python bootstrap.py}, followed by \texttt{python run.py} \\
Documentation & Repository README and \texttt{docs/} \\
Support & GitHub issue tracker; corresponding-author e-mail will be supplied in the journal submission system \\
\bottomrule
\end{tabularx}
\end{table*}

\section{Motivation and significance}\label{sec:motivation}

Sustainability and annual reports are long, visually rich documents in which a single disclosure
may depend on a multi-level table header, a unit outside the selected cell, a footnote, a reporting
boundary and an assurance statement. A fluent answer to a question about such a document is not
therefore equivalent to reusable knowledge. Cross-report analysis additionally requires stable
entity identifiers and qualified predicates so that superficially similar figures are not compared
across incompatible periods, scopes or units.

Existing open-source document converters such as MinerU \citep{wang2024mineru}, Docling
\citep{auer2024docling} and Unstructured \citep{unstructured2026} make PDFs accessible to
downstream language models, but conversion alone does not define when a claim may enter a trusted
knowledge layer. ESG-specific research has used language models to derive structured triples and
graph analyses \citep{bronzini2024glitter}, while ESG-KG integrates multimodal extraction with a
knowledge graph for compliance assessment \citep{chang2026esgkg}. SustainTrace is complementary:
its software contribution is an auditable admission boundary around atomic ESG facts, together
with a separate failure--repair knowledge base and immutable evidence for controlled improvement.
It evaluates fidelity to the supplied report, not the external truth of an issuer's disclosure.

The system produces two linked forms of knowledge. \emph{Report-fact knowledge} represents a
company, predicate, value, period, unit, statistical boundary, method and assurance status together
with source evidence. \emph{Method knowledge} represents the extraction stage, failure signature,
diagnosis, repair strategy, execution state and post-validation outcome. This separation makes a
failed extraction scientifically useful instead of hiding it behind an automatic retry. It also
supports controlled evolution: preserved failures can motivate new components, but an old lockbox
result is never retrospectively changed and a new component must pass on a newly frozen case.

The intended users are sustainability analysts, researchers and software developers who need to
inspect one company, compare qualified facts across reports, or study extraction and repair
failure modes. The release follows research-software citation and FAIR principles
\citep{smith2016software,barker2022fair} and exposes both human-readable interfaces and
machine-readable artifacts.

\section{Software description}\label{sec:software}

\subsection{Architecture and processing contract}

The pipeline has seven auditable stages: (1) freeze a source PDF and register its SHA-256;
(2) parse it with MinerU or a native-PDF/Poppler fallback; (3) project model or rule output into
typed atomic-fact slots; (4) bind values to pages, quotations and available table or spatial
coordinates; (5) apply deterministic integrity, period, unit, boundary, calculation and conflict
gates; (6) route supported candidates to human-gated promotion or archive a typed failure; and
(7) plan and post-validate bounded repairs without modifying the parent record.

Every run records inputs, configuration, hashes, status transitions and outputs. Provenance follows
the same entity--activity--agent separation encouraged by PROV-O \citep{w3cprov2013}, with
additional ESG-specific slots and table coordinates. The trusted fact store and failure--repair
store are queryable together but have different admission policies. A successful parser or
diagnostic component can satisfy only the dimension authorized by its contract; it cannot erase an
unrelated blocker or promote knowledge by itself.

\subsection{Atomic facts, evidence and trust states}

The fact schema covers quantitative indicators, targets and progress, EU-taxonomy disclosures,
reporting boundaries, accounting methods, exclusions, assurance facts, deterministic derivations,
conflicts and explicit absence states. Evidence records identify the source hash, PDF page,
quotation and, where available, the table, row, column, cell and bounding box. These coordinates
allow a user to move from a normalized claim back to the exact disclosure.

Semantic identity and trust are independent. A well-formed record may remain a candidate when one
evidence dimension is unresolved. The release seed places only 149 independently reviewed Tier B
facts in the default trusted layer; candidates, simulated references and failures remain isolated.
Controlled entities have persistent identifiers, aliases and disambiguation rules, and relations
retain period, unit and boundary qualifiers needed for defensible cross-report queries.

\subsection{Refusal and controlled repair}

SustainTrace distinguishes document absence, retrieval failure, parsing failure, semantic mismatch,
unsupported computation and unresolved conflict. The controller maps eligible signatures to
registered repair strategies, checks frozen inputs and executor contracts, and creates a child
candidate with complete lineage. A failed preflight is blocked, a failed post-validation is rolled
back, and a passing repair still requires the applicable promotion policy. Silent retry and
post-hoc semantic completion are prohibited.

\subsection{Workbench and distribution}

The web workbench provides report intake, a knowledge explorer, company analysis, a reliability
center, evaluation summaries and an immutable audit trail. The Knowledge Explorer displays a
readable claim and structured slots beside highlighted source-page evidence. The Reliability
Center exposes failure signatures, repair plans, policies and execution outcomes. The interface
therefore makes both knowledge layers visible rather than presenting only a chat answer.

The same source workflow targets macOS, Windows and Linux. \texttt{bootstrap.py} creates isolated
environments, installs pinned dependencies, verifies model assets, performs a one-page PDF smoke
test, installs the production frontend and initializes the release seed. \texttt{run.py} starts
the API and user interface. Corporate PDFs are not redistributed; users supply or retrieve them
under the publishers' terms.

\section{Illustrative examples}\label{sec:examples}

\subsection{Evidence-linked fact inspection}

An analyst selects a trusted circular-materials fact. The Knowledge Explorer displays its
subject--predicate--value representation, period and unit, then opens the source table on the
registered PDF page with the supporting row highlighted (Fig.~\ref{fig:explorer}). This turns
provenance checking into a direct interface operation rather than a reconstruction from a model
answer.

\begin{figure*}[t]
  \centering
  \submissionfigure[width=\textwidth]{sustaintrace-knowledge-explorer.png}
  \caption{Knowledge Explorer linking a normalized ESG claim to its structured slots and highlighted PDF evidence.}
  \label{fig:explorer}
\end{figure*}

\subsection{Qualified cross-report analysis}

Equivalent indicators can be queried across entities and periods because predicates and entities
are controlled, while units, methods, scopes and reporting boundaries remain attached. A chart or
table can therefore exclude incompatible records instead of silently normalizing them. Explicit
absence and refusal records show where a comparison is incomplete.

\subsection{Failure-to-repair audit}

The Reliability Center groups failures by signature and shows the selected strategy, owner,
execution state and archive status (Fig.~\ref{fig:reliability}). A repair appears as a derived
candidate only after deterministic post-validation; the parent failure and rollback evidence are
preserved. Researchers can inspect which document structures remain difficult and whether a new
component actually generalizes to an unseen frozen case.

\begin{figure*}[t]
  \centering
  \submissionfigure[width=\textwidth]{sustaintrace-reliability-center.png}
  \caption{Reliability Center showing execution states, control policies, repair strategies and archived failure records.}
  \label{fig:reliability}
\end{figure*}

\section{Impact}\label{sec:impact}

\subsection{Evaluation design}

The evaluation asks three questions: whether projected knowledge and evidence match frozen PDFs
(RQ1), whether admission controls prevent unsupported knowledge while retaining useful candidates
(RQ2), and whether repair components and lifecycle controls preserve fail-closed, auditable
behavior (RQ3).

For RQ1 and part of RQ2, two independent annotators reviewed 60 tasks from 20 official reports:
20 direct candidate claims, 20 complex candidates and 20 refusal tasks. An independent adjudicator
resolved 49 task-level and 134 slot-level disagreements. Annotators saw the report and frozen
evidence window but not prompts, simulated references, scores, repair logs, other annotations or
the system decision. Generative-AI assistance was prohibited. The gold data record claim support,
knowledge-base action and slot correctness. Agreement is reported with Cohen's kappa
\citep{cohen1960coefficient}; uncertainty uses 10,000 report-clustered bootstrap repetitions.

The paired E3 experiment isolates the admission gate from generation variability. It contains 120
tasks across 20 reports and six failure strata. Mistral Large 3 and GPT-OSS each generated one
candidate per task. Ungated and full outcomes were derived locally from that identical candidate,
with no second call, silent retry or semantic repair. The source-derived positive and
deterministically mutated negative cases are pseudo-gold admissibility labels, not external factual
truth or a second human benchmark. Discordant decisions use an exact two-sided McNemar test
\citep{mcnemar1947sampling}.

RQ3 uses two non-pooled sets: 13 immutable single-attempt component trials across six component
classes and seven reports, plus four frozen repair-state-machine cases covering promotion,
post-validation rollback and preflight blocking. Component evidence and lifecycle safety are
reported separately because they do not share an accuracy denominator.

\subsection{Knowledge fidelity and admission control}

Raw annotator agreement was 0.567 for the five-level support decision, 0.583 for knowledge-base
action and 0.757 across candidate slots; Cohen's kappa was 0.327, 0.431 and 0.526, respectively.
Among 40 candidate tasks, 20 were fully supported (50.0\%; report-clustered 95\% CI 37.5--65.0\%)
and all were at least partially supported. Applicable slots were exact in 177/220 cases (80.5\%;
95\% CI 73.6--86.7\%), and evidence coordinates were exact in 40/40. Direct claims reached 12/20
exact support versus 8/20 for complex claims. Value was exact in 17/17 applicable cases; subject
identity (22/39) and period binding (20/32) were the largest remaining error sources.

Mistral completed all 120 E3 candidates. Among 100 reference-nonadmissible tasks, unsafe admission
fell from 27/100 ungated to 1/100 with SustainTrace, an absolute reduction of 26.0 percentage points
(95\% CI 17.0--36.0; 26 prevented, none introduced; exact paired $p=2.98\times10^{-8}$). GPT-OSS
completed 119/120; the remaining contract failure was preserved without retry. Among its 99
completed reference-nonadmissible tasks, unsafe admission fell from 24/99 (24.2\%) to 1/99 (1.0\%),
an absolute reduction of 23.2 points (95\% CI 15.2--32.0; 23 prevented, none introduced;
$p=2.38\times10^{-7}$).

The safety gain had a retention cost. Mistral safe admission fell from 12/20 (60.0\%) to 10/20
(50.0\%), retaining 10/12 ungated safe admissions; GPT-OSS fell from 13/20 (65.0\%) to 7/20
(35.0\%), retaining 7/13. The paired effect is therefore consistent in direction, but
model-specific over-blocking requires calibration. Figure~\ref{fig:paired} reports both sides of
this trade-off.

\begin{figure*}[t]
  \centering
  \submissionfigure[width=0.94\textwidth]{sustaintrace-e3-paired-effect-v1.png}
  \caption{Paired admission-control effect on identical frozen candidates with report-clustered 95\% bootstrap intervals: unsafe admission decreases, while safe-admission retention reveals a model-dependent cost.}
  \label{fig:paired}
\end{figure*}

The smaller release audit complements E3. Two of three frozen Tier B promotions were admitted
without correction. No unsupported or unassessable claim entered the trusted layer; one partially
supported claim required bounded correction and is reported separately. The controller was
conservative: 13/20 refusal tasks were human-admissible false abstentions, 14/37 held candidates
were admissible after review, and promotion coverage was 2/16 among all human-admissible candidates.

\subsection{Controlled-repair evidence}

The 13 component trials produced six contract passes and seven preserved failures. Three passes
were positive unseen results for table-cell binding, native-PDF corroboration and calculation
lineage after earlier failures had been retained and used for development. No component trial
wrote a fact or promoted trusted knowledge. The four lifecycle cases ended in one promotion, two
post-validation rollbacks and one preflight block. All reached a safe terminal state; all injected
hash, invariant or preflight faults were detected, and no source or historical lockbox was altered.

These results establish bounded failure-to-new-unseen-pass sequences and lifecycle safety, not
general autonomous repair. No component has repeated unseen success across enough heterogeneous
reports to satisfy the project's highest maturity threshold. Figure~\ref{fig:summary} preserves
the distinct denominators for fidelity, admission and repair evidence.

\begin{figure*}[t]
  \centering
  \submissionfigure[width=0.96\textwidth]{sustaintrace-evaluation-summary-v1.png}
  \caption{Evaluation summary: slot fidelity, admission/abstention outcomes and immutable component history. Failed trials remain visible rather than being rescored.}
  \label{fig:summary}
\end{figure*}

\subsection{Scientific utility and limitations}

SustainTrace makes factual knowledge, provenance and method knowledge about failures available in
one executable system. The release includes schemas, deterministic gates, the web workbench,
machine-readable evaluation artifacts and cross-platform installation. This supports research on
ESG extraction, evidence-based comparison and repair policies without equating fluent generation
with trusted knowledge.

The evidence also defines the limits. Full-claim exact support is 50\%, and safe admission causes
substantial false abstention. Complex tables, subject identity and period binding remain important
error sources. Trusted promotion remains human-gated. E3 demonstrates a paired deterministic gate
effect on pseudo-gold cases; it does not demonstrate external factual truth or fully unseen report
generalization. Some E3 reports occurred in the wider development corpus, and GPT-OSS shows strong
model-specific over-blocking. Finally, the RQ3 evidence does not establish autonomous code
self-modification or universal repair.

\section{Conclusions}\label{sec:conclusions}

SustainTrace is an ESG knowledge infrastructure that converts reports into evidence-linked facts,
supports qualified cross-report analysis and accumulates auditable failure--repair knowledge for
controlled evolution. Its contribution is not merely a new answer generator: it defines when a
claim may enter a trusted layer, preserves refusals and failures, and validates bounded repairs
without rewriting historical evidence. Human-adjudicated and paired evaluations support evidence
traceability and fail-closed admission while quantifying current fidelity and retention limits.
The open software and immutable artifacts provide a reproducible basis for improving ESG document
intelligence without overstating factual truth or autonomous repair.

\section*{Data and code availability}
Source code is available at \url{https://github.com/extradimen/sustaintrace} under Apache-2.0 and
is archived at \url{https://doi.org/10.5281/zenodo.23293034}.
Derived records and evaluation artifacts are distributable with the software. Issuer PDFs remain
under their original copyright and are not redistributed.

\section*{Ethics statement}
Ethics approval was not required. The evaluation involved professional annotation of public
corporate reports and software outputs, collected no sensitive personal data, and involved no
intervention with human participants.

\section*{Declaration of competing interest}
The authors declare that they have no known competing financial interests or personal
relationships that could have appeared to influence the work reported in this paper.

\section*{Acknowledgements}
This research received no specific grant from any funding agency in the public, commercial, or
not-for-profit sectors.

\printcredits

\begin{thebibliography}{99}

\bibitem[Wang et al.(2024)]{wang2024mineru}
B. Wang, C. Xu, X. Zhao, et al.,
MinerU: An open-source solution for precise document content extraction,
arXiv:2409.18839 (2024).
\newblock \url{https://doi.org/10.48550/arXiv.2409.18839}.

\bibitem[Auer et al.(2024)]{auer2024docling}
C. Auer, M. Lysak, A. Nassar, et al.,
Docling technical report,
arXiv:2408.09869 (2024).
\newblock \url{https://doi.org/10.48550/arXiv.2408.09869}.

\bibitem[Unstructured Technologies(2026)]{unstructured2026}
Unstructured Technologies,
Unstructured: Open-source preprocessing tools for unstructured data (2026).
\newblock \url{https://docs.unstructured.io/}.

\bibitem[Bronzini et al.(2024)]{bronzini2024glitter}
M. Bronzini, C. Nicolini, B. Lepri, A. Passerini, J. Staiano,
Glitter or gold? Deriving structured insights from sustainability reports via large language models,
EPJ Data Science 13 (2024) 41.
\newblock \url{https://doi.org/10.1140/epjds/s13688-024-00481-2}.

\bibitem[Chang et al.(2026)]{chang2026esgkg}
L.-Y. Chang, C.-M. Chen, H.-H. Huang, M.-F. Tsai, A.-Z. Yen, C.-J. Wang,
ESG-KG: A multi-modal knowledge graph system for automated compliance assessment,
in: Proceedings of EACL System Demonstrations, 2026, pp. 602--608.
\newblock \url{https://doi.org/10.18653/v1/2026.eacl-demo.43}.

\bibitem[Smith et al.(2016)]{smith2016software}
A.M. Smith, D.S. Katz, K.E. Niemeyer, FORCE11 Software Citation Working Group,
Software citation principles,
PeerJ Computer Science 2 (2016) e86.
\newblock \url{https://doi.org/10.7717/peerj-cs.86}.

\bibitem[Barker et al.(2022)]{barker2022fair}
M. Barker, N.P. Chue Hong, D.S. Katz, et al.,
Introducing the FAIR principles for research software,
Scientific Data 9 (2022) 622.
\newblock \url{https://doi.org/10.1038/s41597-022-01710-x}.

\bibitem[W3C Provenance Working Group(2013)]{w3cprov2013}
W3C Provenance Working Group,
PROV-O: The PROV ontology, W3C Recommendation (2013).
\newblock \url{https://www.w3.org/TR/prov-o/}.

\bibitem[Cohen(1960)]{cohen1960coefficient}
J. Cohen,
A coefficient of agreement for nominal scales,
Educational and Psychological Measurement 20 (1960) 37--46.
\newblock \url{https://doi.org/10.1177/001316446002000104}.

\bibitem[McNemar(1947)]{mcnemar1947sampling}
Q. McNemar,
Note on the sampling error of the difference between correlated proportions or percentages,
Psychometrika 12 (1947) 153--157.
\newblock \url{https://doi.org/10.1007/BF02295996}.

\end{thebibliography}

\end{document}
